\section{Objet et périmètre du code}
\label{sec-objet}

\subsection{Le problème physique}

Le forage percussif fragmente la roche par une succession de chocs : un piston frappe un taillant (le bit), l'onde de compression traverse le bit et les inserts en carbure de tungstène la transmettent à la roche, qui se fissure, se broie et éjecte des fragments sous chaque insert. rockim simule ce problème à l'échelle d'un insert : l'élasticité du volume, la fissuration par des joints cohésifs, le contact frottant entre fragments, l'effet de la vitesse de déformation sur la résistance, la pulvérisation du volume et le bilan d'énergie de tout le système (\cle{LISEZ_MOI.md}, l.~6-13).

Deux familles de données servent de cibles, et elles ne portent pas sur la même roche. Les essais d'impact à insert unique d'Imperial College et de Mines Paris-PSL, publiés par Yang et al.~\cite{yang2025validation,yang2026pulverisation}, portent sur le calcaire de St Anne, le grès de Rhune et le granite de Kuru Grey (\cle{specs/005-impact-insert-yang/spec.md}, l.~9-24). Les essais quasi statiques sur le granite de Red Bohus, matériau de la thèse, fournissent une résistance en compression simple de $126{,}6 \pm 21{,}4$~MPa, une résistance brésilienne de $10{,}27$~MPa, un module de $77{,}7$~GPa et un coefficient de Poisson de $0{,}29$~\cite{dumoulin2024data} (\cle{calibration_redbohus/README.md}, l.~18-25). Aucune donnée d'impact sur le Red Bohus n'alimente aujourd'hui rockim.

\subsection{Raison d'être d'un code maison}

Le dépôt ne contient pas de justification écrite du choix de développer un code plutôt que d'utiliser un code existant. Les documents fixent quatre objectifs au code : une parité structurelle avec le code de l'article de référence (MultiFracS, lignée Munjiza), la reproduction d'un impact validable contre le banc de Mines Paris, la robustesse et la vitesse (\cle{ROADMAP_rockim.md}, l.~3-6). Le guide d'utilisation le présente comme un bac à sable d'exploration, de calibration et de figures, la production 3D calibrée restant confiée à Abaqus et aux VUMAT (\cle{GUIDE_rockim.md}, l.~465-466).

Le code de référence des articles de Yang et al., Solidity (Imperial College, licence LGPL version 3), n'est pas redistribué dans rockim. rockim en transcrit des conventions, ligne à ligne pour certaines, en citant les fichiers sources de Solidity. Le dépôt public de rockim n'a pas de licence alors qu'il contient ces transcriptions : ce point juridique est ouvert (\cle{HANDOFF_2026-10-03.md}, l.~192-193). Une contrainte du doctorant distingue aussi rockim de Solidity : garder l'insertion adaptative des joints (Yan et al. 2023~\cite{yan2023adaptive}), là où Solidity place les joints partout dès l'origine.

\subsection{Les modes de calcul}

Six solveurs partagent une interface commune et se choisissent par la clé \cle{mode} (\cref{tab-modes}). Chaque mode accepte plusieurs scénarios de chargement : \cle{percussion} (impacteur libre), \cle{shear} (coupe par un outil traîné à profondeur imposée, dont un couteau PDC en 2D et en 3D), \cle{tension} (traction, ou compression simple et triaxiale par platines et confinement), \cle{loads} (charges par groupes de nœuds, ajouté le 3 octobre 2026), \cle{brazilian} et \cle{shpb} en FDEM 2D seulement, et \cle{jointbench} (banc d'un joint isolé) en FDEM 3D. Douze sous-commandes \cle{selftest-*} vérifient des briques sans maillage (\cref{sec-verification}).

\begin{table}[htbp]
\centering
\caption{Les six solveurs de rockim et leur statut depuis la décision de périmètre du 14 août 2026.}
\label{tab-modes}
\begin{tabularx}{\linewidth}{@{}l X l l@{}}
\toprule
mode & modèle & rupture & statut \\
\midrule
\cle{fem} & éléments finis 2D, déformation plane, triangles à déformation constante & endommagement et érosion & gelé \\
\cle{fem3d} & tétraèdres co-rotationnels, lois de volume au choix & endommagement et érosion & gelé, enrichi \\
\cle{dem}, \cle{dem3d} & particules liées (modèle à liaisons parallèles) & rupture des liaisons & gelé \\
\cle{fdem} & FDEM 2D, grains de Voronoï possibles & joints cohésifs explicites & actif \\
\cle{fdem3d} & FDEM 3D, tétraèdres, grains de Voronoï possibles & joints cohésifs triangulaires & actif \\
\bottomrule
\end{tabularx}
\end{table}

La décision du 14 août 2026 concentre le développement sur \cle{fdem} et \cle{fdem3d} (\cle{ROADMAP_rockim.md}, l.~10-29). Les modes gelés gardent leurs auto-tests. Le mode \cle{fem3d} a pourtant reçu après cette date une loi d'endommagement plastique (le 4 septembre), un matériau par phase (le 6 septembre) et les charges par groupes (le 3 octobre) : il sert de contrepartie à Abaqus sur maillage identique et de référence élastique dans la campagne de vérification.

Trois usages servent réellement la thèse. Le mode \cle{fdem3d} porte les répliques d'impact de Yang et al. (St Anne et Kuru). Le mode \cle{fdem} en 2D porte la calibration sur le Red Bohus, avec 279 jeux de données d'essai (189 en traction ou compression, 85 brésiliens, 5 de coupe), calibration à refaire depuis la correction des amortissements (\cle{ROADMAP_rockim.md}, l.~207). Le mode \cle{fem3d} sert à la validation croisée avec Abaqus.

\subsection{Architecture}

Le code source compte 18 fichiers dans \cle{src/} et 28 en-têtes dans \cle{include/rockim/}, soit 47\,302 lignes au 6 octobre 2026. Les plus gros fichiers sont les solveurs FDEM 2D (10\,957 lignes) et 3D (8\,817 lignes) et la bibliothèque de lois de volume \cle{MatLaw.cpp} (5\,511 lignes). Des briques transverses portent la lecture des configurations (\cle{Config}), le registre des clés admises par mode (\cle{KeyGuard.hpp}, \cle{KeysByMode.hpp}), la géométrie du contact par potentiel (\cle{PotentialContact.hpp}), les lois cohésives (\cle{JointTsl.hpp}), les charges par groupes et l'écriture des fichiers VTK. Les en-têtes portent le raisonnement physique et la provenance des choix (\cle{LISEZ_MOI.md}, l.~85).

Une constitution de huit principes encadre le développement (\cle{.specify/memory/constitution.md}). Trois d'entre eux gouvernent la lecture des résultats de ce rapport. Une option nouvelle est désactivée par défaut et son absence laisse les sorties identiques au bit près, ce que vérifie un jeu d'ancres d'empreintes. Le bilan d'énergie est le juge physique de tout changement. Le code croît par addition : une loi remplacée reste disponible derrière sa clé.

\subsection{Chronologie}

L'historique git commence le 19 août 2026 ; l'antérieur n'est connu que par les documents, et la date de naissance de la première version (FEM et DEM 2D) n'y figure pas. Les documents datent les étapes suivantes.
\begin{itemize}
\item 7 août 2026 : instabilité 3D, $2{,}4$~MJ créés pour un impact de 16~J, due au calcul du pas de temps sur des tétraèdres de Kuhn.
\item 11 au 14 août : audit complet, documentation extraite du code, activation adaptative du contact, contact par potentiel en 2D et en 3D, corps maillés nommés, bilan d'énergie à sept postes, décision de périmètre.
\item 19 au 22 août : couplage hydromécanique 2D, confronté au cas d'AbuAisha et al.~\cite{abuaisha2017hm}.
\item 5 septembre : arbre de développement g0, compilation CMake, ancre d'identité au bit sur huit configurations, refus des clés inconnues.
\item 11 au 14 septembre : campagne de reproduction des articles de Yang et al., transcription de la loi de joint de Solidity, contact par potentiel parallélisé, calcul St Anne de $300~\mu$s.
\item 3 octobre : charges par groupes, campagne d'optimisation, pas de temps variable, force de contact par volume de recouvrement.
\item 3 et 4 octobre : campagne de vérification et validation, bancs de l'onde dans une barre et du bloc glissant.
\end{itemize}
Les 109 commits postérieurs au 25 août vivent sur une branche qui n'est pas fusionnée dans la branche principale du dépôt et n'a été ni compilée ni testée sous Windows (\cle{HANDOFF_2026-10-03.md}, l.~9-10 et 29-30).
